\documentclass[12pt,a4paper]{article}

\usepackage[utf8]{inputenc}
\usepackage[T2A]{fontenc}
\usepackage[russian]{babel}
\usepackage{amsmath,amssymb,amsthm}
\usepackage{array}
\usepackage{booktabs}
\usepackage{enumitem}
\usepackage{listings}
\usepackage{xcolor}
\usepackage{geometry}

\geometry{margin=2.2cm}

\newtheorem{definition}{Определение}[section]
\newtheorem{theorem}{Теорема}[section]
\newtheorem{lemma}{Лемма}[section]
\newtheorem{proposition}{Утверждение}[section]
\newtheorem{example}{Пример}[section]

\theoremstyle{remark}
\newtheorem{remark}{Замечание}[section]

\lstdefinestyle{sqlstyle}{
    language=SQL,
    basicstyle=\ttfamily\small,
    keywordstyle=\color{blue},
    commentstyle=\color{gray},
    stringstyle=\color{red!60!black},
    breaklines=true,
    frame=single
}

\title{Билет №72 \\ \small Теоретические основы реляционной модели данных (РДМ). Реляционная алгебра, реляционное исчисление. Функциональные зависимости и нормализация отношений.(ИТМО) \\ Основные понятия реляционной и объектной моделей данных. (СпБГУ)}
\author{Сериков Владислав Александрович}
\date{16 мая 2026}

\begin{document}

\maketitle
\tableofcontents
\newpage

\section{Модель данных: общее понятие}

\begin{definition}
\textbf{Модель данных} --- это совокупность понятий, правил и ограничений, предназначенных для описания структуры данных, допустимых операций над ними и условий целостности.
\end{definition}

Обычно в модели данных выделяют три компонента:
\begin{enumerate}
    \item \textbf{структурная часть} --- какие объекты используются для представления данных;
    \item \textbf{операционная часть} --- какие операции допустимы над данными;
    \item \textbf{ограничения целостности} --- какие состояния базы данных считаются корректными.
\end{enumerate}

Классическими моделями данных являются:
\begin{itemize}
    \item иерархическая модель;
    \item сетевая модель;
    \item реляционная модель;
    \item объектная модель;
    \item объектно-реляционная модель;
    \item полуструктурированные и документо-ориентированные модели.
\end{itemize}

Главное внимание в данном конспекте уделяется реляционной и объектной моделям.

\section{Теоретические основы реляционной модели данных}

\subsection{Основная идея реляционной модели}

Реляционная модель данных была предложена Э. Ф. Коддом как математически строгая модель представления данных на основе понятия отношения.

Основные идеи реляционной модели:
\begin{enumerate}
    \item данные представляются в виде отношений, то есть таблиц;
    \item строки таблицы соответствуют кортежам;
    \item столбцы таблицы соответствуют атрибутам;
    \item порядок строк и порядок столбцов несущественен;
    \item операции над данными формально описываются средствами реляционной алгебры или реляционного исчисления;
    \item целостность данных задаётся ограничениями.
\end{enumerate}

\subsection{Домены, атрибуты, кортежи и отношения}

\begin{definition}
\textbf{Домен} --- это множество допустимых атомарных значений некоторого типа.
\end{definition}

Например:
\[
D_{\text{StudentId}} = \{1,2,3,\dots\},
\]
\[
D_{\text{Name}} = \{\text{строки длиной не более 100 символов}\}.
\]

\begin{definition}
\textbf{Атрибут} --- это именованная роль домена в отношении.
\end{definition}

Например, атрибуты:
\[
\text{StudentId}, \quad \text{Name}, \quad \text{GroupNo}.
\]

\begin{definition}
Пусть заданы атрибуты
\[
A_1, A_2, \dots, A_n,
\]
каждый из которых имеет домен
\[
D_1, D_2, \dots, D_n.
\]
Тогда \textbf{кортеж} над этими атрибутами --- это функция
\[
t : \{A_1,\dots,A_n\} \to D_1 \cup \dots \cup D_n,
\]
такая что
\[
t(A_i) \in D_i.
\]
\end{definition}

В табличной форме кортеж соответствует строке.

\begin{definition}
\textbf{Схема отношения} --- это выражение вида
\[
R(A_1, A_2, \dots, A_n),
\]
где $R$ --- имя отношения, а $A_1,\dots,A_n$ --- имена атрибутов.
\end{definition}

\begin{definition}
\textbf{Отношение} $r$ со схемой $R(A_1,\dots,A_n)$ --- это конечное множество кортежей над атрибутами $A_1,\dots,A_n$.
\end{definition}

Если атрибуты имеют домены $D_1,\dots,D_n$, то отношение можно рассматривать как конечное подмножество декартова произведения:
\[
r \subseteq D_1 \times D_2 \times \dots \times D_n.
\]

\begin{example}
Рассмотрим отношение:
\[
\text{Student}(\text{sid}, \text{name}, \text{group})
\]

\[
\begin{array}{c|c|c}
\text{sid} & \text{name} & \text{group} \\
\hline
1 & \text{Иванов} & \text{M3435} \\
2 & \text{Петров} & \text{M3435} \\
3 & \text{Сидоров} & \text{M3436}
\end{array}
\]

Здесь:
\begin{itemize}
    \item \texttt{Student} --- имя отношения;
    \item \texttt{sid}, \texttt{name}, \texttt{group} --- атрибуты;
    \item каждая строка --- кортеж;
    \item множество всех строк --- отношение.
\end{itemize}
\end{example}

\subsection{Свойства отношений}

В реляционной модели отношение обладает следующими свойствами:
\begin{enumerate}
    \item отсутствуют дублирующиеся кортежи, поскольку отношение является множеством;
    \item порядок кортежей несущественен;
    \item порядок атрибутов несущественен, если атрибуты имеют имена;
    \item значения атрибутов атомарны;
    \item все значения одного атрибута принадлежат одному домену.
\end{enumerate}

\subsection{Схема базы данных и состояние базы данных}

\begin{definition}
\textbf{Схема базы данных} --- это множество схем отношений:
\[
\mathcal{R} = \{R_1, R_2, \dots, R_m\}.
\]
\end{definition}

\begin{definition}
\textbf{Состояние базы данных} --- это множество конкретных отношений:
\[
I = \{r_1, r_2, \dots, r_m\},
\]
где каждое $r_i$ является отношением со схемой $R_i$.
\end{definition}

\subsection{Ключи}

\begin{definition}
Пусть $R$ --- схема отношения, $r$ --- отношение со схемой $R$, а $X \subseteq R$ --- множество атрибутов. Говорят, что $X$ является \textbf{суперключом} отношения $R$, если для любых двух различных кортежей $t_1,t_2 \in r$ выполняется:
\[
t_1[X] \neq t_2[X].
\]
\end{definition}

Иначе говоря, значения атрибутов из $X$ однозначно идентифицируют кортеж.

\begin{definition}
\textbf{Ключ-кандидат} --- это минимальный по включению суперключ.
\end{definition}

\begin{definition}
\textbf{Первичный ключ} --- это один из ключей-кандидатов, выбранный в качестве основного идентификатора кортежей.
\end{definition}

\begin{definition}
\textbf{Внешний ключ} --- это множество атрибутов одного отношения, значения которого должны совпадать со значениями ключа в другом отношении.
\end{definition}

\begin{example}
Пусть имеются отношения:
\[
\text{Student}(\text{sid}, \text{name}, \text{group\_id})
\]
\[
\text{Group}(\text{group\_id}, \text{title})
\]

Тогда:
\begin{itemize}
    \item \texttt{Student.sid} может быть первичным ключом отношения \texttt{Student};
    \item \texttt{Group.group\_id} может быть первичным ключом отношения \texttt{Group};
    \item \texttt{Student.group\_id} является внешним ключом, ссылающимся на \texttt{Group.group\_id}.
\end{itemize}
\end{example}

\subsection{Ограничения целостности}

В реляционной модели обычно выделяют следующие ограничения:
\begin{enumerate}
    \item \textbf{доменная целостность}: каждое значение принадлежит соответствующему домену;
    \item \textbf{целостность сущностей}: первичный ключ не должен содержать неопределённых значений;
    \item \textbf{ссылочная целостность}: значение внешнего ключа должно ссылаться на существующий кортеж или быть допустимо неопределённым;
    \item \textbf{пользовательские ограничения}: дополнительные условия предметной области.
\end{enumerate}

\section{Реляционная алгебра}

\subsection{Понятие реляционной алгебры}

\begin{definition}
\textbf{Реляционная алгебра} --- это формальный процедурный язык запросов, операции которого принимают на вход отношения и возвращают отношения.
\end{definition}

Основное свойство реляционной алгебры --- \textbf{замкнутость}: результат любой операции снова является отношением.

Это позволяет строить сложные выражения из простых операций.

\subsection{Базовые операции реляционной алгебры}

К классическим операциям реляционной алгебры относятся:
\begin{enumerate}
    \item выборка;
    \item проекция;
    \item объединение;
    \item разность;
    \item декартово произведение;
    \item переименование.
\end{enumerate}

На их основе определяются:
\begin{enumerate}
    \item пересечение;
    \item соединение;
    \item естественное соединение;
    \item деление;
    \item полусоединение;
    \item внешние соединения.
\end{enumerate}

\subsection{Выборка}

\begin{definition}
\textbf{Выборка} $\sigma_\theta(r)$ --- это операция, возвращающая все кортежи отношения $r$, удовлетворяющие условию $\theta$:
\[
\sigma_\theta(r) = \{t \in r \mid \theta(t) = \text{true}\}.
\]
\end{definition}

\begin{example}
Пусть:
\[
\text{Student} =
\begin{array}{c|c|c}
\text{sid} & \text{name} & \text{group} \\
\hline
1 & \text{Иванов} & \text{M3435} \\
2 & \text{Петров} & \text{M3435} \\
3 & \text{Сидоров} & \text{M3436}
\end{array}
\]

Запрос:
\[
\sigma_{\text{group} = \text{``M3435''}}(\text{Student})
\]

Результат:
\[
\begin{array}{c|c|c}
\text{sid} & \text{name} & \text{group} \\
\hline
1 & \text{Иванов} & \text{M3435} \\
2 & \text{Петров} & \text{M3435}
\end{array}
\]
\end{example}

\subsection{Проекция}

\begin{definition}
\textbf{Проекция} $\pi_{A_1,\dots,A_k}(r)$ --- это операция, возвращающая отношение, состоящее только из указанных атрибутов:
\[
\pi_{A_1,\dots,A_k}(r) =
\{t[A_1,\dots,A_k] \mid t \in r\}.
\]
\end{definition}

Так как отношение является множеством, дубликаты после проекции удаляются.

\begin{example}
\[
\pi_{\text{group}}(\text{Student})
\]

Результат:
\[
\begin{array}{c}
\text{group} \\
\hline
\text{M3435} \\
\text{M3436}
\end{array}
\]
\end{example}

\subsection{Объединение}

\begin{definition}
Пусть отношения $r$ и $s$ имеют одинаковые схемы. Тогда их \textbf{объединение}:
\[
r \cup s = \{t \mid t \in r \vee t \in s\}.
\]
\end{definition}

Операция объединения определена только для совместимых по объединению отношений.

\begin{definition}
Отношения $r$ и $s$ называются \textbf{совместимыми по объединению}, если:
\begin{enumerate}
    \item они имеют одинаковую степень;
    \item соответствующие атрибуты имеют одинаковые домены.
\end{enumerate}
\end{definition}

\subsection{Разность}

\begin{definition}
Пусть $r$ и $s$ совместимы по объединению. Тогда:
\[
r - s = \{t \mid t \in r \wedge t \notin s\}.
\]
\end{definition}

\subsection{Декартово произведение}

\begin{definition}
Пусть $r(R)$ и $s(S)$ --- отношения, причём множества атрибутов $R$ и $S$ не пересекаются. Тогда \textbf{декартово произведение}:
\[
r \times s = \{tq \mid t \in r,\ q \in s\},
\]
где $tq$ --- конкатенация кортежей $t$ и $q$.
\end{definition}

Если имена атрибутов совпадают, предварительно используется переименование.

\subsection{Переименование}

\begin{definition}
Операция \textbf{переименования} обозначается:
\[
\rho_{S(B_1,\dots,B_n)}(r)
\]
и возвращает отношение с новым именем $S$ и, возможно, новыми именами атрибутов $B_1,\dots,B_n$.
\end{definition}

\subsection{Производные операции}

\subsubsection{Пересечение}

\begin{definition}
Пересечение двух совместимых отношений:
\[
r \cap s = \{t \mid t \in r \wedge t \in s\}.
\]
\end{definition}

Через базовые операции:
\[
r \cap s = r - (r - s).
\]

\subsubsection{Соединение}

\begin{definition}
\textbf{$\theta$-соединение} отношений $r$ и $s$:
\[
r \bowtie_\theta s = \sigma_\theta(r \times s).
\]
\end{definition}

\begin{example}
Пусть:
\[
\text{Student}(\text{sid}, \text{name}, \text{group\_id})
\]
\[
\text{Group}(\text{group\_id}, \text{title})
\]

Тогда запрос:
\[
\text{Student} \bowtie_{\text{Student.group\_id} =
\text{Group.group\_id}} \text{Group}
\]
соединяет студентов с их группами.
\end{example}

\subsubsection{Естественное соединение}

\begin{definition}
\textbf{Естественное соединение} $r \bowtie s$ соединяет кортежи по всем одноимённым атрибутам и оставляет один экземпляр каждого общего атрибута.
\end{definition}

Если:
\[
R = X \cup Y,\quad S = Y \cup Z,
\]
то:
\[
r \bowtie s =
\pi_{X,Y,Z}
\left(
\sigma_{r.Y = s.Y}(r \times s)
\right).
\]

\subsubsection{Деление}

\begin{definition}
Пусть $r(X,Y)$ и $s(Y)$ --- отношения. \textbf{Деление} $r \div s$ возвращает множество таких значений $X$, которые связаны со всеми значениями $Y$ из $s$:
\[
r \div s =
\{x \mid \forall y \in s: (x,y) \in r\}.
\]
\end{definition}

Через базовые операции:
\[
r \div s =
\pi_X(r) -
\pi_X \left(
(\pi_X(r) \times s) - r
\right).
\]

\begin{example}
Пусть отношение \texttt{Passed(student, course)} показывает, какие курсы сдал студент:
\[
\begin{array}{c|c}
\text{student} & \text{course} \\
\hline
\text{Иванов} & \text{БД} \\
\text{Иванов} & \text{Алгоритмы} \\
\text{Петров} & \text{БД} \\
\text{Сидоров} & \text{БД} \\
\text{Сидоров} & \text{Алгоритмы}
\end{array}
\]

Пусть:
\[
\text{Required} =
\begin{array}{c}
\text{course} \\
\hline
\text{БД} \\
\text{Алгоритмы}
\end{array}
\]

Тогда:
\[
\text{Passed} \div \text{Required}
\]
вернёт студентов, сдавших все обязательные курсы:
\[
\begin{array}{c}
\text{student} \\
\hline
\text{Иванов} \\
\text{Сидоров}
\end{array}
\]
\end{example}

\subsection{Эквивалентность выражений реляционной алгебры}

Некоторые полезные эквивалентности:

\[
\sigma_{\theta_1 \wedge \theta_2}(r)
=
\sigma_{\theta_1}(\sigma_{\theta_2}(r)).
\]

\[
\pi_X(\pi_Y(r)) = \pi_X(r),
\quad \text{если } X \subseteq Y.
\]

\[
\sigma_\theta(r \bowtie s)
=
\sigma_\theta(r) \bowtie s,
\]
если условие $\theta$ использует только атрибуты отношения $r$.

\[
r \bowtie s = s \bowtie r.
\]

\[
(r \bowtie s) \bowtie q = r \bowtie (s \bowtie q).
\]

Эти эквивалентности применяются при оптимизации запросов.

\section{Реляционное исчисление}

\subsection{Общая характеристика}

\begin{definition}
\textbf{Реляционное исчисление} --- это декларативный формальный язык запросов, основанный на логике предикатов первого порядка.
\end{definition}

В отличие от реляционной алгебры, реляционное исчисление описывает, \emph{что} нужно получить, а не \emph{как} это вычислять.

Выделяют:
\begin{enumerate}
    \item исчисление кортежей;
    \item исчисление доменов.
\end{enumerate}

\subsection{Исчисление кортежей}

\begin{definition}
В \textbf{исчислении кортежей} переменные пробегают кортежи отношений.
Запрос имеет вид:
\[
\{t \mid P(t)\},
\]
где $P(t)$ --- формула логики первого порядка.
\end{definition}

Формулы строятся из:
\begin{itemize}
    \item атомов вида $R(t)$, означающих, что кортеж $t$ принадлежит отношению $R$;
    \item сравнений $t.A \ \theta \ c$;
    \item сравнений $t.A \ \theta \ s.B$;
    \item логических связок $\wedge, \vee, \neg$;
    \item кванторов $\exists, \forall$.
\end{itemize}

\begin{example}
Найти студентов группы M3435:
\[
\{t \mid \text{Student}(t) \wedge t.\text{group} = \text{``M3435''}\}.
\]
\end{example}

\begin{example}
Найти имена студентов, сдавших курс ``БД'':
\[
\{t.\text{name} \mid
\text{Student}(t) \wedge
\exists p(
\text{Passed}(p)
\wedge p.\text{sid}=t.\text{sid}
\wedge p.\text{course}=\text{``БД''})
\}.
\]
\end{example}

\subsection{Исчисление доменов}

\begin{definition}
В \textbf{исчислении доменов} переменные пробегают значения доменов.
Запрос имеет вид:
\[
\{ \langle x_1,\dots,x_n\rangle \mid P(x_1,\dots,x_n)\}.
\]
\end{definition}

\begin{example}
Найти имена студентов группы M3435:
\[
\{ \langle n\rangle \mid
\exists id \exists g(
\text{Student}(id,n,g)
\wedge g = \text{``M3435''})
\}.
\]
\end{example}

\subsection{Безопасность выражений}

Не всякая формула реляционного исчисления задаёт конечный результат.

\begin{example}
Формула:
\[
\{x \mid \neg \text{Student}(x)\}
\]
потенциально задаёт бесконечное множество всех кортежей, не являющихся студентами.
\end{example}

\begin{definition}
Выражение реляционного исчисления называется \textbf{безопасным}, если его результат конечен и содержит только значения, встречающиеся в базе данных или в самом запросе.
\end{definition}

\subsection{Теорема Кодда об эквивалентности}

\begin{theorem}[Теорема Кодда]
Реляционная алгебра и безопасное реляционное исчисление имеют одинаковую выразительную силу: для каждого выражения реляционной алгебры существует эквивалентное безопасное выражение реляционного исчисления, и для каждого безопасного выражения реляционного исчисления существует эквивалентное выражение реляционной алгебры.
\end{theorem}

\begin{proof}
Доказательство проводится в две стороны.

\textbf{1. Из реляционной алгебры в реляционное исчисление.}

Покажем индукцией по построению выражения реляционной алгебры, что каждому выражению можно сопоставить формулу исчисления.

База индукции: если выражение является именем отношения $R$, то ему соответствует формула:
\[
R(t).
\]

Переход индукции:
\begin{enumerate}
    \item Выборка:
    \[
    \sigma_\theta(E)
    \]
    переводится в:
    \[
    \{t \mid \varphi_E(t) \wedge \theta(t)\}.
    \]

    \item Проекция:
    если
    \[
    E' = \pi_{A_1,\dots,A_k}(E),
    \]
    то:
    \[
    \{u \mid \exists t(\varphi_E(t)
    \wedge u.A_1=t.A_1 \wedge \dots \wedge u.A_k=t.A_k)\}.
    \]

    \item Объединение:
    \[
    E_1 \cup E_2
    \]
    переводится в:
    \[
    \{t \mid \varphi_{E_1}(t) \vee \varphi_{E_2}(t)\}.
    \]

    \item Разность:
    \[
    E_1 - E_2
    \]
    переводится в:
    \[
    \{t \mid \varphi_{E_1}(t) \wedge \neg \varphi_{E_2}(t)\}.
    \]

    \item Декартово произведение:
    \[
    E_1 \times E_2
    \]
    переводится в формулу, утверждающую существование кортежей $t_1$ и $t_2$, принадлежащих результатам $E_1$ и $E_2$, таких что результирующий кортеж является их конкатенацией.
\end{enumerate}

Таким образом, любая операция реляционной алгебры переводится в формулу исчисления.

\textbf{2. Из безопасного реляционного исчисления в реляционную алгебру.}

Формулы безопасного реляционного исчисления строятся из атомов, логических связок и кванторов. Для каждой такой конструкции можно построить соответствующее выражение алгебры:
\begin{enumerate}
    \item атом $R(t)$ соответствует отношению $R$;
    \item конъюнкция соответствует соединению или выборке;
    \item дизъюнкция соответствует объединению;
    \item отрицание в безопасных формулах соответствует разности относительно конечной активной области;
    \item существующий квантор соответствует проекции;
    \item универсальный квантор выражается через отрицание и существование:
    \[
    \forall x\,\varphi(x) \equiv \neg \exists x\, \neg \varphi(x).
    \]
\end{enumerate}

Безопасность необходима, чтобы отрицание и кванторы рассматривались только над конечной активной областью базы данных. Следовательно, результат можно представить конечным выражением реляционной алгебры.

Обе стороны доказаны, значит реляционная алгебра и безопасное реляционное исчисление эквивалентны по выразительной силе.
\end{proof}

\section{Функциональные зависимости}

\subsection{Понятие функциональной зависимости}

\begin{definition}
Пусть $R$ --- схема отношения, а $X,Y \subseteq R$ --- множества атрибутов.
Говорят, что на $R$ выполняется \textbf{функциональная зависимость}
\[
X \to Y,
\]
если для любых двух кортежей $t_1,t_2$ любого допустимого отношения $r(R)$ из равенства
\[
t_1[X] = t_2[X]
\]
следует
\[
t_1[Y] = t_2[Y].
\]
\end{definition}

Иначе говоря, значения атрибутов $X$ однозначно определяют значения атрибутов $Y$.

\begin{example}
В отношении:
\[
\text{Student}(\text{sid}, \text{name}, \text{group})
\]
обычно выполняется зависимость:
\[
\text{sid} \to \text{name}, \text{group}.
\]
Это означает, что идентификатор студента однозначно определяет имя и группу.
\end{example}

\subsection{Тривиальные и нетривиальные зависимости}

\begin{definition}
Функциональная зависимость $X \to Y$ называется \textbf{тривиальной}, если
\[
Y \subseteq X.
\]
\end{definition}

Например:
\[
\{\text{sid}, \text{name}\} \to \{\text{name}\}
\]
является тривиальной.

\begin{definition}
Функциональная зависимость $X \to Y$ называется \textbf{нетривиальной}, если
\[
Y \nsubseteq X.
\]
\end{definition}

\subsection{Замыкание множества функциональных зависимостей}

\begin{definition}
Пусть $F$ --- множество функциональных зависимостей. \textbf{Замыкание} $F^+$ --- это множество всех функциональных зависимостей, логически следующих из $F$.
\end{definition}

То есть:
\[
F^+ = \{X \to Y \mid F \models X \to Y\}.
\]

\subsection{Аксиомы Армстронга}

\begin{theorem}[Аксиомы Армстронга]
Для функциональных зависимостей корректны следующие правила вывода:
\begin{enumerate}
    \item \textbf{рефлексивность}: если $Y \subseteq X$, то
    \[
    X \to Y;
    \]
    \item \textbf{пополнение}: если
    \[
    X \to Y,
    \]
    то для любого $Z$:
    \[
    XZ \to YZ;
    \]
    \item \textbf{транзитивность}: если
    \[
    X \to Y
    \quad \text{и} \quad
    Y \to Z,
    \]
    то
    \[
    X \to Z.
    \]
\end{enumerate}
\end{theorem}

\begin{proof}
Докажем корректность каждого правила.

\textbf{1. Рефлексивность.}

Пусть $Y \subseteq X$. Возьмём любые два кортежа $t_1,t_2$ такие, что:
\[
t_1[X] = t_2[X].
\]
Так как $Y$ является подмножеством $X$, равенство значений на $X$ влечёт равенство значений на $Y$:
\[
t_1[Y] = t_2[Y].
\]
Следовательно:
\[
X \to Y.
\]

\textbf{2. Пополнение.}

Пусть $X \to Y$. Нужно доказать:
\[
XZ \to YZ.
\]
Возьмём любые кортежи $t_1,t_2$ такие, что:
\[
t_1[XZ] = t_2[XZ].
\]
Тогда одновременно:
\[
t_1[X] = t_2[X]
\]
и
\[
t_1[Z] = t_2[Z].
\]
Из $X \to Y$ следует:
\[
t_1[Y] = t_2[Y].
\]
Следовательно:
\[
t_1[YZ] = t_2[YZ].
\]
Значит:
\[
XZ \to YZ.
\]

\textbf{3. Транзитивность.}

Пусть:
\[
X \to Y,
\quad
Y \to Z.
\]
Возьмём любые кортежи $t_1,t_2$ такие, что:
\[
t_1[X] = t_2[X].
\]
Из зависимости $X \to Y$ получаем:
\[
t_1[Y] = t_2[Y].
\]
Из зависимости $Y \to Z$ получаем:
\[
t_1[Z] = t_2[Z].
\]
Следовательно:
\[
X \to Z.
\]
\end{proof}

\subsection{Производные правила вывода}

Из аксиом Армстронга выводятся следующие правила:

\begin{enumerate}
    \item \textbf{объединение}:
    \[
    X \to Y,\ X \to Z \Rightarrow X \to YZ;
    \]

    \item \textbf{декомпозиция}:
    \[
    X \to YZ \Rightarrow X \to Y,\ X \to Z;
    \]

    \item \textbf{псевдотранзитивность}:
    \[
    X \to Y,\ WY \to Z \Rightarrow WX \to Z.
    \]
\end{enumerate}

\begin{proof}
Докажем правило объединения.

Пусть:
\[
X \to Y
\quad \text{и} \quad
X \to Z.
\]
По правилу пополнения из $X \to Y$ получаем:
\[
XX \to XY.
\]
Так как $XX=X$, имеем:
\[
X \to XY.
\]
Из $X \to Z$ по пополнению с $Y$:
\[
XY \to ZY.
\]
По транзитивности из
\[
X \to XY
\quad \text{и} \quad
XY \to ZY
\]
получаем:
\[
X \to ZY.
\]
То есть:
\[
X \to YZ.
\]

Правило декомпозиции следует из рефлексивности и транзитивности:
если $X \to YZ$, а $Y \subseteq YZ$, то $YZ \to Y$, значит $X \to Y$.
Аналогично $X \to Z$.

Правило псевдотранзитивности:
из $X \to Y$ по пополнению с $W$ получаем:
\[
WX \to WY.
\]
Из $WY \to Z$ и транзитивности следует:
\[
WX \to Z.
\]
\end{proof}

\subsection{Замыкание множества атрибутов}

\begin{definition}
Пусть $F$ --- множество функциональных зависимостей, а $X$ --- множество атрибутов.
\textbf{Замыкание} $X^+$ относительно $F$ --- это множество всех атрибутов, функционально зависящих от $X$:
\[
X^+ = \{A \mid F \models X \to A\}.
\]
\end{definition}

\subsection{Алгоритм вычисления замыкания множества атрибутов}

\textbf{Вход:}
\begin{itemize}
    \item множество атрибутов $X$;
    \item множество функциональных зависимостей $F$.
\end{itemize}

\textbf{Выход:}
\[
X^+.
\]

\textbf{Алгоритм:}
\begin{enumerate}
    \item Положить:
    \[
    X^+ := X.
    \]
    \item Пока существуют зависимости $Y \to Z$ из $F$, такие что:
    \[
    Y \subseteq X^+
    \quad \text{и} \quad
    Z \nsubseteq X^+,
    \]
    добавить атрибуты $Z$ в $X^+$:
    \[
    X^+ := X^+ \cup Z.
    \]
    \item Когда новых атрибутов добавить нельзя, вернуть $X^+$.
\end{enumerate}

\begin{example}
Пусть:
\[
R(A,B,C,D,E),
\]
\[
F = \{A \to B,\ B \to C,\ AC \to D,\ D \to E\}.
\]

Найдём $A^+$.

\[
A^+ := \{A\}.
\]

Так как $A \to B$, добавляем $B$:
\[
A^+ = \{A,B\}.
\]

Так как $B \to C$, добавляем $C$:
\[
A^+ = \{A,B,C\}.
\]

Так как $AC \to D$ и $A,C \in A^+$, добавляем $D$:
\[
A^+ = \{A,B,C,D\}.
\]

Так как $D \to E$, добавляем $E$:
\[
A^+ = \{A,B,C,D,E\}.
\]

Следовательно:
\[
A^+ = R.
\]
Значит, $A$ является суперключом.
\end{example}

\subsection{Проверка функциональной зависимости}

Чтобы проверить, следует ли зависимость $X \to Y$ из $F$, достаточно вычислить $X^+$ и проверить:
\[
Y \subseteq X^+.
\]

\begin{proposition}
\[
F \models X \to Y
\quad \Longleftrightarrow \quad
Y \subseteq X^+.
\]
\end{proposition}

\begin{proof}
По определению:
\[
X^+ = \{A \mid F \models X \to A\}.
\]

Если $Y \subseteq X^+$, то каждый атрибут из $Y$ функционально зависит от $X$. По правилу объединения:
\[
X \to Y.
\]

Если $F \models X \to Y$, то по правилу декомпозиции для каждого $A \in Y$:
\[
F \models X \to A,
\]
значит:
\[
A \in X^+.
\]
Следовательно:
\[
Y \subseteq X^+.
\]
\end{proof}

\subsection{Минимальное покрытие функциональных зависимостей}

\begin{definition}
\textbf{Минимальное покрытие} множества функциональных зависимостей $F$ --- это такое множество зависимостей $G$, что:
\[
F^+ = G^+,
\]
и выполняются условия:
\begin{enumerate}
    \item правая часть каждой зависимости содержит один атрибут;
    \item в левой части каждой зависимости нет лишних атрибутов;
    \item в $G$ нет лишних зависимостей.
\end{enumerate}
\end{definition}

\subsection{Алгоритм построения минимального покрытия}

\textbf{Вход:}
множество функциональных зависимостей $F$.

\textbf{Выход:}
минимальное покрытие $G$.

\textbf{Алгоритм:}
\begin{enumerate}
    \item Разбить зависимости с несколькими атрибутами справа:
    \[
    X \to A_1A_2\dots A_k
    \]
    заменить на:
    \[
    X \to A_1,\ X \to A_2,\dots,\ X \to A_k.
    \]

    \item Для каждой зависимости $X \to A$ и каждого атрибута $B \in X$ проверить, является ли $B$ лишним:
    \[
    (X - \{B\})^+ \supseteq \{A\}.
    \]
    Если да, заменить $X \to A$ на:
    \[
    (X - \{B\}) \to A.
    \]

    \item Для каждой зависимости $f \in G$ проверить, следует ли она из $G-\{f\}$.
    Если да, удалить $f$.
\end{enumerate}

\begin{example}
Пусть:
\[
F = \{A \to BC,\ B \to C,\ A \to B,\ AB \to C\}.
\]

Шаг 1. Разбиваем правые части:
\[
G = \{A \to B,\ A \to C,\ B \to C,\ A \to B,\ AB \to C\}.
\]

Удалим дубликат:
\[
G = \{A \to B,\ A \to C,\ B \to C,\ AB \to C\}.
\]

Шаг 2. Проверим лишние атрибуты в $AB \to C$.

Проверим, лишний ли $A$:
\[
B^+ = \{B,C\},
\]
значит $B \to C$, поэтому $A$ лишний в $AB \to C$.

Заменяем:
\[
AB \to C
\]
на:
\[
B \to C.
\]

Получаем:
\[
G = \{A \to B,\ A \to C,\ B \to C\}.
\]

Шаг 3. Проверяем лишнюю зависимость $A \to C$.

Из:
\[
A \to B,\quad B \to C
\]
по транзитивности следует:
\[
A \to C.
\]

Поэтому удаляем $A \to C$.

Итог:
\[
G = \{A \to B,\ B \to C\}.
\]
\end{example}

\section{Нормализация отношений}

\subsection{Проблема избыточности}

Нормализация предназначена для устранения нежелательных свойств схем отношений:
\begin{enumerate}
    \item избыточного хранения данных;
    \item аномалий вставки;
    \item аномалий удаления;
    \item аномалий обновления.
\end{enumerate}

\begin{example}
Пусть имеется отношение:
\[
\text{StudentCourse}(\text{sid}, \text{name}, \text{course}, \text{teacher})
\]

\[
\begin{array}{c|c|c|c}
\text{sid} & \text{name} & \text{course} & \text{teacher} \\
\hline
1 & \text{Иванов} & \text{БД} & \text{Смирнов} \\
1 & \text{Иванов} & \text{Алгоритмы} & \text{Петров} \\
2 & \text{Сидоров} & \text{БД} & \text{Смирнов}
\end{array}
\]

Функциональные зависимости:
\[
\text{sid} \to \text{name},
\]
\[
\text{course} \to \text{teacher}.
\]

Проблемы:
\begin{itemize}
    \item имя студента повторяется для каждого курса;
    \item преподаватель курса повторяется для каждого студента;
    \item при изменении преподавателя курса нужно менять много строк;
    \item нельзя добавить курс без студента, если \texttt{sid} входит в ключ.
\end{itemize}
\end{example}

\subsection{Декомпозиция схем}

\begin{definition}
\textbf{Декомпозиция} схемы отношения $R$ --- это замена $R$ набором схем:
\[
R_1,\dots,R_k,
\]
таких что:
\[
R = R_1 \cup R_2 \cup \dots \cup R_k.
\]
\end{definition}

Декомпозиция должна по возможности обладать двумя важными свойствами:
\begin{enumerate}
    \item соединение без потерь;
    \item сохранение зависимостей.
\end{enumerate}

\subsection{Соединение без потерь}

\begin{definition}
Декомпозиция $R$ на $R_1$ и $R_2$ называется \textbf{декомпозицией с соединением без потерь} относительно множества функциональных зависимостей $F$, если для любого допустимого отношения $r(R)$:
\[
r = \pi_{R_1}(r) \bowtie \pi_{R_2}(r).
\]
\end{definition}

\begin{theorem}[Критерий соединения без потерь для бинарной декомпозиции]
Пусть схема $R$ декомпозирована на $R_1$ и $R_2$.
Декомпозиция является декомпозицией с соединением без потерь относительно $F$ тогда и только тогда, когда:
\[
(R_1 \cap R_2) \to R_1
\]
или
\[
(R_1 \cap R_2) \to R_2
\]
следует из $F^+$.
\end{theorem}

\begin{proof}
Обозначим:
\[
X = R_1 \cap R_2.
\]

\textbf{Достаточность.}

Пусть:
\[
X \to R_1.
\]
Нужно доказать:
\[
r = \pi_{R_1}(r) \bowtie \pi_{R_2}(r).
\]

Всегда верно:
\[
r \subseteq \pi_{R_1}(r) \bowtie \pi_{R_2}(r),
\]
так как каждый кортеж исходного отношения после проекций соединится сам с собой.

Докажем обратное включение.

Пусть $t$ --- кортеж из:
\[
\pi_{R_1}(r) \bowtie \pi_{R_2}(r).
\]
Тогда существуют кортежи $u,v \in r$ такие, что:
\[
t[R_1] = u[R_1],
\]
\[
t[R_2] = v[R_2],
\]
и они совпадают на общих атрибутах $X$:
\[
u[X] = v[X].
\]

Из зависимости:
\[
X \to R_1
\]
следует:
\[
u[R_1] = v[R_1].
\]

Тогда кортеж $v$ совпадает с $t$ на $R_1$ и на $R_2$, следовательно:
\[
t = v.
\]
Поскольку $v \in r$, получаем:
\[
t \in r.
\]

Значит:
\[
\pi_{R_1}(r) \bowtie \pi_{R_2}(r) \subseteq r.
\]

Следовательно:
\[
r = \pi_{R_1}(r) \bowtie \pi_{R_2}(r).
\]

Случай $X \to R_2$ доказывается аналогично.

\textbf{Необходимость.}

Покажем идею доказательства.

Если ни
\[
X \to R_1,
\]
ни
\[
X \to R_2
\]
не следуют из $F^+$, то существуют допустимые отношения, в которых два кортежа могут совпадать на $X$, но различаться в атрибутах $R_1-X$ и $R_2-X$.

Тогда после проекции на $R_1$ и $R_2$ естественное соединение создаст дополнительный ложный кортеж, которого не было в исходном отношении. Следовательно, декомпозиция не будет без потерь.

Значит, для соединения без потерь необходимо, чтобы общая часть функционально определяла одну из компонент.
\end{proof}

\begin{example}
Пусть:
\[
R(A,B,C),
\quad F=\{A \to B\}.
\]

Декомпозиция:
\[
R_1(A,B),
\quad R_2(A,C).
\]

Здесь:
\[
R_1 \cap R_2 = \{A\}.
\]

Так как:
\[
A \to B,
\]
то:
\[
A \to AB = R_1.
\]

Следовательно, декомпозиция без потерь.
\end{example}

\subsection{Сохранение зависимостей}

\begin{definition}
Декомпозиция $R$ на $R_1,\dots,R_k$ сохраняет зависимости $F$, если:
\[
\left(
\pi_{R_1}(F) \cup \dots \cup \pi_{R_k}(F)
\right)^+
=
F^+.
\]
\end{definition}

Здесь $\pi_{R_i}(F)$ --- множество функциональных зависимостей из $F^+$, содержащих только атрибуты $R_i$.

Практически сохранение зависимостей означает, что все ограничения можно проверять локально, не соединяя отношения.

\subsection{Первая нормальная форма}

\begin{definition}
Отношение находится в \textbf{первой нормальной форме} (1НФ), если значения всех его атрибутов атомарны.
\end{definition}

Нарушение 1НФ:
\[
\begin{array}{c|c}
\text{student} & \text{phones} \\
\hline
\text{Иванов} & \{111,222\}
\end{array}
\]

Приведение к 1НФ:
\[
\begin{array}{c|c}
\text{student} & \text{phone} \\
\hline
\text{Иванов} & 111 \\
\text{Иванов} & 222
\end{array}
\]

\subsection{Вторая нормальная форма}

\begin{definition}
Неключевой атрибут --- это атрибут, не входящий ни в один ключ-кандидат.
\end{definition}

\begin{definition}
Отношение находится во \textbf{второй нормальной форме} (2НФ), если оно находится в 1НФ и каждый неключевой атрибут полностью функционально зависит от каждого ключа-кандидата.
\end{definition}

\begin{definition}
Функциональная зависимость $X \to A$ называется \textbf{частичной зависимостью}, если существует собственное подмножество $Y \subset X$ такое, что:
\[
Y \to A.
\]
\end{definition}

\begin{example}
Пусть:
\[
R(\text{student}, \text{course}, \text{student\_name}, \text{grade})
\]
и ключ:
\[
(\text{student}, \text{course}).
\]

Зависимости:
\[
\text{student} \to \text{student\_name},
\]
\[
(\text{student}, \text{course}) \to \text{grade}.
\]

Атрибут \texttt{student\_name} зависит только от части составного ключа. Это нарушение 2НФ.

Декомпозиция:
\[
R_1(\text{student}, \text{student\_name}),
\]
\[
R_2(\text{student}, \text{course}, \text{grade}).
\]
\end{example}

\subsection{Третья нормальная форма}

\begin{definition}
Пусть $R$ --- схема отношения, $F$ --- множество функциональных зависимостей.
Отношение $R$ находится в \textbf{третьей нормальной форме} (3НФ), если для любой нетривиальной функциональной зависимости
\[
X \to A
\]
из $F^+$ выполняется хотя бы одно из условий:
\begin{enumerate}
    \item $X$ является суперключом;
    \item $A$ является простым атрибутом, то есть входит хотя бы в один ключ-кандидат.
\end{enumerate}
\end{definition}

\begin{example}
Пусть:
\[
R(\text{student}, \text{group}, \text{curator})
\]
и зависимости:
\[
\text{student} \to \text{group},
\]
\[
\text{group} \to \text{curator}.
\]

Ключ:
\[
\text{student}.
\]

Зависимость:
\[
\text{group} \to \text{curator}
\]
нарушает 3НФ, поскольку \texttt{group} не является суперключом, а \texttt{curator} не является простым атрибутом.

Декомпозиция:
\[
R_1(\text{student}, \text{group}),
\]
\[
R_2(\text{group}, \text{curator}).
\]
\end{example}

\subsection{Нормальная форма Бойса--Кодда}

\begin{definition}
Отношение $R$ находится в \textbf{нормальной форме Бойса--Кодда} (НФБК, BCNF), если для любой нетривиальной функциональной зависимости
\[
X \to Y
\]
из $F^+$ множество $X$ является суперключом.
\end{definition}

BCNF строже, чем 3НФ.

\begin{proposition}
Если отношение находится в BCNF, то оно находится в 3НФ.
\end{proposition}

\begin{proof}
Пусть отношение находится в BCNF. Тогда для любой нетривиальной зависимости:
\[
X \to A
\]
левая часть $X$ является суперключом.

Но условие 3НФ требует, чтобы для любой такой зависимости выполнялось хотя бы одно из условий:
\begin{enumerate}
    \item $X$ --- суперключ;
    \item $A$ --- простой атрибут.
\end{enumerate}

Первое условие всегда выполняется. Следовательно, отношение находится в 3НФ.
\end{proof}

\subsection{Пример различия 3НФ и BCNF}

Пусть:
\[
R(A,B,C)
\]
и зависимости:
\[
F = \{AB \to C,\ C \to B\}.
\]

Ключи:
\[
AB,\quad AC.
\]

Атрибуты $A,B,C$ являются простыми, так как каждый входит хотя бы в один ключ.

Зависимость:
\[
C \to B
\]
не нарушает 3НФ, потому что $B$ --- простой атрибут.

Но $C$ не является суперключом, поэтому зависимость $C \to B$ нарушает BCNF.

Следовательно, отношение находится в 3НФ, но не находится в BCNF.

\subsection{Алгоритм декомпозиции в BCNF}

\textbf{Вход:}
схема $R$ и множество функциональных зависимостей $F$.

\textbf{Выход:}
декомпозиция $R_1,\dots,R_k$ в BCNF с соединением без потерь.

\textbf{Алгоритм:}
\begin{enumerate}
    \item Если все текущие схемы находятся в BCNF, завершить.
    \item Найти схему $S$, не находящуюся в BCNF.
    \item Найти нетривиальную зависимость:
    \[
    X \to Y,
    \]
    нарушающую BCNF, то есть $X$ не является суперключом схемы $S$.
    \item Заменить $S$ двумя схемами:
    \[
    S_1 = X \cup Y,
    \]
    \[
    S_2 = S - (Y - X).
    \]
    \item Повторять до тех пор, пока все схемы не окажутся в BCNF.
\end{enumerate}

\begin{example}
Пусть:
\[
R(A,B,C),
\quad F=\{A \to B,\ B \to C\}.
\]

Ключ:
\[
A,
\]
так как:
\[
A^+ = \{A,B,C\}.
\]

Зависимость:
\[
B \to C
\]
нарушает BCNF, поскольку $B$ не является суперключом.

Декомпозиция:
\[
R_1(B,C),
\]
\[
R_2(A,B).
\]

Проверим:
\begin{itemize}
    \item в $R_1(B,C)$ зависимость $B \to C$, и $B$ --- ключ;
    \item в $R_2(A,B)$ зависимость $A \to B$, и $A$ --- ключ.
\end{itemize}

Обе схемы находятся в BCNF.

Декомпозиция без потерь, поскольку:
\[
R_1 \cap R_2 = \{B\},
\]
а:
\[
B \to BC = R_1.
\]
\end{example}

\subsection{Алгоритм синтеза схемы в 3НФ}

\textbf{Вход:}
схема $R$ и множество функциональных зависимостей $F$.

\textbf{Выход:}
декомпозиция в 3НФ, сохраняющая зависимости и имеющая соединение без потерь.

\textbf{Алгоритм:}
\begin{enumerate}
    \item Найти минимальное покрытие $G$ множества $F$.
    \item Для каждой зависимости:
    \[
    X \to A
    \]
    из $G$ создать схему:
    \[
    R_i = X \cup \{A\}.
    \]
    \item Если ни одна из полученных схем не содержит ключ исходной схемы $R$, добавить отдельную схему, содержащую некоторый ключ $R$.
    \item Удалить схемы, являющиеся подмножествами других схем.
\end{enumerate}

\begin{example}
Пусть:
\[
R(A,B,C,D),
\]
\[
F=\{A \to B,\ B \to C,\ C \to D\}.
\]

Минимальное покрытие уже имеет вид:
\[
G=\{A \to B,\ B \to C,\ C \to D\}.
\]

Создаём схемы:
\[
R_1(A,B),
\quad
R_2(B,C),
\quad
R_3(C,D).
\]

Ключ исходной схемы:
\[
A,
\]
так как:
\[
A^+ = \{A,B,C,D\}.
\]

Схема $R_1(A,B)$ содержит ключ $A$, поэтому дополнительная схема не нужна.

Итоговая декомпозиция:
\[
R_1(A,B),\ R_2(B,C),\ R_3(C,D).
\]
\end{example}

\subsection{Четвёртая нормальная форма}

\begin{definition}
Многозначная зависимость
\[
X \twoheadrightarrow Y
\]
в схеме $R$ означает, что для каждого значения $X$ множество значений $Y$ не зависит от множества значений остальных атрибутов:
\[
R - X - Y.
\]
\end{definition}

\begin{definition}
Отношение находится в \textbf{четвёртой нормальной форме} (4НФ), если для каждой нетривиальной многозначной зависимости:
\[
X \twoheadrightarrow Y
\]
множество $X$ является суперключом.
\end{definition}

\begin{example}
Пусть:
\[
R(\text{student}, \text{language}, \text{hobby}).
\]

Если языки студента и его хобби независимы, то:
\[
\text{student} \twoheadrightarrow \text{language},
\]
\[
\text{student} \twoheadrightarrow \text{hobby}.
\]

Для студента с двумя языками и двумя хобби появится четыре строки, что создаёт избыточность.

Декомпозиция:
\[
R_1(\text{student}, \text{language}),
\]
\[
R_2(\text{student}, \text{hobby}).
\]
\end{example}

\section{Основные понятия объектной модели данных}

\subsection{Общая идея объектной модели}

\begin{definition}
\textbf{Объектная модель данных} --- это модель данных, в которой данные представляются в виде объектов, обладающих состоянием, поведением и идентичностью.
\end{definition}

Объектная модель возникла под влиянием объектно-ориентированного программирования и предназначена для более естественного представления сложных структур данных.

В объектной модели основными понятиями являются:
\begin{enumerate}
    \item объект;
    \item идентификатор объекта;
    \item состояние;
    \item методы;
    \item класс;
    \item тип;
    \item наследование;
    \item инкапсуляция;
    \item полиморфизм;
    \item ассоциации между объектами;
    \item сложные объекты и коллекции.
\end{enumerate}

\subsection{Объект}

\begin{definition}
\textbf{Объект} --- это сущность, имеющая:
\begin{enumerate}
    \item уникальную идентичность;
    \item состояние;
    \item поведение.
\end{enumerate}
\end{definition}

Например, объект \texttt{student1} может иметь состояние:
\[
\text{name} = \text{``Иванов''},
\quad
\text{group} = \text{``M3435''}.
\]

И поведение:
\[
\texttt{enroll(course)},
\quad
\texttt{changeGroup(newGroup)}.
\]

\subsection{Идентичность объекта}

\begin{definition}
\textbf{Идентификатор объекта} --- это системно поддерживаемая сущность, позволяющая отличать объект от всех остальных объектов независимо от значений его атрибутов.
\end{definition}

В реляционной модели идентификация обычно осуществляется через значения ключей. В объектной модели два объекта могут иметь одинаковые значения полей, но оставаться различными объектами.

\begin{example}
Два студента могут иметь одинаковые имя и дату рождения, но быть разными объектами:
\[
o_1 \neq o_2,
\]
хотя:
\[
o_1.\text{name} = o_2.\text{name},
\quad
o_1.\text{birthDate} = o_2.\text{birthDate}.
\]
\end{example}

\subsection{Состояние объекта}

\begin{definition}
\textbf{Состояние объекта} --- это совокупность значений его атрибутов в данный момент времени.
\end{definition}

Атрибуты объекта могут быть:
\begin{itemize}
    \item атомарными;
    \item ссылочными;
    \item структурными;
    \item коллекционными.
\end{itemize}

\begin{example}
Объект \texttt{Student} может содержать:
\begin{itemize}
    \item \texttt{name : String};
    \item \texttt{age : Integer};
    \item \texttt{group : Group};
    \item \texttt{courses : Set<Course>}.
\end{itemize}
\end{example}

\subsection{Методы и поведение}

\begin{definition}
\textbf{Метод} --- это операция, связанная с объектом или классом и определяющая поведение объекта.
\end{definition}

Например:
\[
\texttt{Student.enroll(Course c)}
\]
добавляет курс в множество курсов студента.

В отличие от реляционной модели, где операции над данными обычно задаются внешним языком запросов, объектная модель связывает данные и операции в одном объекте.

\subsection{Классы и типы}

\begin{definition}
\textbf{Класс} --- это описание множества объектов с одинаковой структурой и поведением.
\end{definition}

\begin{definition}
\textbf{Тип} --- это множество допустимых значений и операций над ними.
\end{definition}

В объектных системах класс часто одновременно играет роль:
\begin{itemize}
    \item шаблона для создания объектов;
    \item типа переменных;
    \item контейнера методов.
\end{itemize}

\begin{example}
\begin{lstlisting}[style=sqlstyle]
class Student {
    String name;
    Group group;
    Set<Course> courses;

    void enroll(Course c) {
        courses.add(c);
    }
}
\end{lstlisting}
\end{example}

\subsection{Инкапсуляция}

\begin{definition}
\textbf{Инкапсуляция} --- это принцип, согласно которому внутреннее состояние объекта скрывается, а доступ к нему осуществляется через методы.
\end{definition}

Например, вместо прямого изменения поля:
\[
student.group = newGroup
\]
используется метод:
\[
student.changeGroup(newGroup).
\]

Это позволяет контролировать корректность изменений.

\subsection{Наследование}

\begin{definition}
\textbf{Наследование} --- это отношение между классами, при котором один класс получает атрибуты и методы другого класса.
\end{definition}

Если класс \texttt{Student} наследуется от класса \texttt{Person}, то он получает атрибуты:
\[
\texttt{name},\quad \texttt{birthDate}
\]
и добавляет собственные:
\[
\texttt{group},\quad \texttt{studentId}.
\]

\begin{example}
\begin{lstlisting}[style=sqlstyle]
class Person {
    String name;
    Date birthDate;
}

class Student extends Person {
    String group;
    Integer studentId;
}
\end{lstlisting}
\end{example}

В базах данных наследование может использоваться для моделирования иерархий сущностей:
\[
\text{Person}
\]
\[
\quad \uparrow
\]
\[
\text{Student},\quad \text{Teacher}.
\]

\subsection{Полиморфизм}

\begin{definition}
\textbf{Полиморфизм} --- это возможность использовать объекты разных классов единообразно, если они имеют общий тип или общий интерфейс.
\end{definition}

Например, если \texttt{Student} и \texttt{Teacher} наследуются от \texttt{Person}, то коллекция:
\[
\texttt{Set<Person>}
\]
может содержать и студентов, и преподавателей.

\subsection{Сложные объекты}

\begin{definition}
\textbf{Сложный объект} --- это объект, состояние которого включает другие объекты или коллекции объектов.
\end{definition}

Пример:
\[
\texttt{Student}
\]
может содержать коллекцию:
\[
\texttt{Set<Course>}.
\]

Объектная модель естественно поддерживает:
\begin{itemize}
    \item вложенные структуры;
    \item списки;
    \item множества;
    \item массивы;
    \item ссылки между объектами;
    \item графы объектов.
\end{itemize}

\subsection{Связи между объектами}

Связи могут быть:
\begin{enumerate}
    \item \textbf{ассоциациями};
    \item \textbf{агрегациями};
    \item \textbf{композициями};
    \item \textbf{наследованием}.
\end{enumerate}

\begin{definition}
\textbf{Ассоциация} --- это логическая связь между объектами.
\end{definition}

Например:
\[
\texttt{Student} \leftrightarrow \texttt{Course}.
\]

\begin{definition}
\textbf{Агрегация} --- это связь ``целое--часть'', при которой часть может существовать независимо от целого.
\end{definition}

\begin{definition}
\textbf{Композиция} --- это строгая форма агрегации, при которой часть не существует независимо от целого.
\end{definition}

\section{Сравнение реляционной и объектной моделей данных}

\subsection{Основные различия}

\[
\begin{array}{p{5cm}|p{5cm}|p{5cm}}
\toprule
\textbf{Критерий} & \textbf{Реляционная модель} & \textbf{Объектная модель} \\
\midrule
Базовая единица & Отношение, кортеж & Объект \\
\midrule
Идентификация & Ключи & Идентификаторы объектов \\
\midrule
Структура & Плоские таблицы & Сложные вложенные структуры \\
\midrule
Операции & Реляционная алгебра, SQL & Методы объектов, объектные запросы \\
\midrule
Связи & Внешние ключи, соединения & Ссылки между объектами \\
\midrule
Наследование & Отсутствует в классической модели & Поддерживается естественно \\
\midrule
Инкапсуляция & Обычно отсутствует & Является базовым принципом \\
\midrule
Математическая основа & Теория множеств и логика предикатов & Объектно-ориентированная парадигма \\
\bottomrule
\end{array}
\]

\subsection{Достоинства реляционной модели}

\begin{enumerate}
    \item строгая математическая основа;
    \item простота табличного представления;
    \item декларативный язык запросов SQL;
    \item независимость данных от физического хранения;
    \item развитая теория нормализации;
    \item эффективные методы оптимизации запросов;
    \item широкая промышленная поддержка.
\end{enumerate}

\subsection{Недостатки реляционной модели}

\begin{enumerate}
    \item неудобство представления сложных объектов;
    \item необходимость соединений для восстановления объектных структур;
    \item отсутствие встроенных методов и поведения;
    \item проблема объектно-реляционного несоответствия при работе с ООП.
\end{enumerate}

\subsection{Достоинства объектной модели}

\begin{enumerate}
    \item естественное представление сложных объектов;
    \item поддержка идентичности объектов;
    \item поддержка инкапсуляции и методов;
    \item поддержка наследования и полиморфизма;
    \item удобство для инженерных, графовых, мультимедийных и CAD/CAM-приложений.
\end{enumerate}

\subsection{Недостатки объектной модели}

\begin{enumerate}
    \item более сложная формальная теория;
    \item менее универсальный язык запросов по сравнению с SQL;
    \item сложность оптимизации запросов;
    \item меньшая стандартизация и распространённость;
    \item трудности обеспечения независимости данных от программ.
\end{enumerate}

\section{Объектно-реляционная модель}

\subsection{Мотивация}

Объектно-реляционная модель объединяет идеи реляционной и объектной моделей.

Она сохраняет табличное представление и SQL, но добавляет:
\begin{itemize}
    \item пользовательские типы;
    \item структурированные типы;
    \item коллекции;
    \item наследование типов или таблиц;
    \item методы;
    \item ссылки;
    \item расширяемость системы типов.
\end{itemize}

\subsection{Пример объектно-реляционного подхода}

В объектно-реляционной СУБД можно определить составной тип:

\begin{lstlisting}[style=sqlstyle]
CREATE TYPE address_type AS (
    city text,
    street text,
    house text
);
\end{lstlisting}

И использовать его в таблице:

\begin{lstlisting}[style=sqlstyle]
CREATE TABLE students (
    sid integer PRIMARY KEY,
    name text,
    address address_type
);
\end{lstlisting}

Такое представление нарушает классическую атомарность первой нормальной формы, но может быть удобным для прикладной разработки.

\subsection{Отношение объектно-реляционной модели к классической РДМ}

Объектно-реляционная модель не отменяет реляционную модель, а расширяет её. Однако расширения могут ослаблять строгость классической реляционной теории, особенно если допускаются:
\begin{itemize}
    \item вложенные отношения;
    \item неатомарные значения;
    \item ссылки вместо внешних ключей;
    \item методы с побочными эффектами.
\end{itemize}

\section{Связь реляционной алгебры с SQL}

SQL является промышленным языком реляционных и объектно-реляционных СУБД, но не полностью совпадает с классической реляционной алгеброй.

Основные отличия SQL от классической реляционной модели:
\begin{enumerate}
    \item SQL допускает дубликаты строк;
    \item SQL использует неопределённое значение \texttt{NULL};
    \item SQL основан на трёхзначной логике;
    \item порядок строк может задаваться через \texttt{ORDER BY};
    \item поддерживаются агрегатные функции;
    \item поддерживаются внешние соединения;
    \item современные реализации поддерживают объектно-реляционные расширения.
\end{enumerate}

Пример соответствия операций:

\[
\sigma_{\text{group}=\text{``M3435''}}(\text{Student})
\]

соответствует SQL-запросу:

\begin{lstlisting}[style=sqlstyle]
SELECT *
FROM Student
WHERE group_no = 'M3435';
\end{lstlisting}

Проекция:
\[
\pi_{\text{name}}(\text{Student})
\]

соответствует:

\begin{lstlisting}[style=sqlstyle]
SELECT DISTINCT name
FROM Student;
\end{lstlisting}

Естественное соединение:

\begin{lstlisting}[style=sqlstyle]
SELECT *
FROM Student
NATURAL JOIN Groups;
\end{lstlisting}

\section{Итоговая схема основных понятий}

\[
\begin{array}{c}
\text{Модель данных} \\
\downarrow \\
\text{Реляционная модель} \\
\downarrow \\
\text{Отношение, атрибут, домен, кортеж} \\
\downarrow \\
\text{Ключи и ограничения целостности} \\
\downarrow \\
\text{Реляционная алгебра и реляционное исчисление} \\
\downarrow \\
\text{Функциональные зависимости} \\
\downarrow \\
\text{Нормализация: 1НФ, 2НФ, 3НФ, BCNF, 4НФ}
\end{array}
\]

\[
\begin{array}{c}
\text{Объектная модель} \\
\downarrow \\
\text{Объект} \\
\downarrow \\
\text{Идентичность, состояние, поведение} \\
\downarrow \\
\text{Классы, типы, методы} \\
\downarrow \\
\text{Инкапсуляция, наследование, полиморфизм} \\
\downarrow \\
\text{Сложные объекты и связи}
\end{array}
\]

\end{document}
